\documentclass[10pt,addpoints,answers]{exam}
\usepackage{graphicx} % Required for inserting images
\usepackage{booktabs}
\graphicspath{{bilag/}}


\begin{document}
\begin{center}
{\Huge \textbf{Kravanalyse + (Test)}}\\[0.5em]
\end{center}

\section*{Krav}




\section*{FURPS+ oversigt over funktionelle krav (F) og ikke-funktionelle krav (URPS+)}
\begin{center}
\renewcommand{\arraystretch}{1.4}
\begin{tabular}{| p{50pt} | p{400pt} |}
\hline
\textbf{FURPS+} & \textbf{Krav} \\ \hline
F1 & Mikrofoner kan detektere en lydhændelse over 75 decibel, målt fra 1 meter \\ \hline
F2 & Systemet kan lokalisere lyden ved hjælp af mikrofonerne, indenfor 30 cm \\ \hline
F3 & Systemet skal pege en markør der peger på målet, med en præcision på 40 cm  \\ \hline
F4 & Systemet har et måleområde på 3x3 meter, hvor mikrofoner placeres i hvert hjørne  \\ \hline
F5 & Tui viser hver hændelse i 1 min før det slettes. Derved kan flere hændelser vises på samme tid\\ \hline
F6 & Systemet implementeres med 2 mikrofon nodes, og en orchestrator node. Mikrofon nodes er synkroniseret via PTP \\ \hline
U1 & TUI til bruger, hvor bruger kan se lokation af hændelsen, samt indstille systemet. \\ \hline
R2 & Sensormodulerne skal kunne genstarte ved crash og gendanne forbindelse til PTP, samt streaming netværket igen automatisk \\ \hline
R3 & Systemet skal kunne identificerer over 90\% af hændelserne. \\ \hline
R4 & Ved usikre estimeringer (residualvector > xx), fortæller systemet at estimeringen er usikker.\\ \hline
P2 & Sensormodulernes clock skal være synkroniseret med maksimalt 200 mikrosekunder \\ \hline
P3 & Position vises senest 1 s efter hændelsen, i TUI. \\ \hline
S1 & Mikrofon antal og positioner skal kunne indstilles i runtime via TUI \\
\hline
\end{tabular}
\renewcommand{\arraystretch}{1}
\end{center}

\newpage
\section*{Use cases}

\subsection*{UC1 - Lokaliser lydhændelse}

\begin{center}
\renewcommand{\arraystretch}{1.4}
\begin{tabular}{| p{150pt} | p{300pt} |}
\hline

\textbf{Navn:} & Lokaliser lydhændelse \\ \hline
\textbf{Mål:} & Operatøren får vist den estimerede position af en lydhændelse, der opstår i måleområdet \\ \hline
\textbf{Initiering:} & En lydhændelse opstår i måleområdet \\ \hline
\textbf{Aktører:} & Operatør (primær), Lydkilde (sekundær, passiv) \\ \hline
\textbf{Antal samtidige forekomster:} & 1 \\ \hline
\textbf{Prækondition:} & UC2 konfigurér opstilling, sensormoduler online og synkroniserede \\ \hline
\textbf{Postkondition:} & Lydhændelsens estimerede position og tidspunkt er vist for operatøren. \\ \hline
\textbf{Hovedscenarie:} & 
1. En lydhændelse opstår i måleområdet. \newline
2. Alle systemets mikrofoner detekterer lydhændelsen. \newline
3. Systemet bestemmer ankomsttidspunktet for hver mikrofon. \newline 
4. Systemet estimerer lydhændelsens position ud fra forskellene i ankomsttid.\newline 
5. Systemet vurderer at estimatet er troværdigt (ikke-funktionelt krav R6). \newline
6. Systemet viser den estimerede position og tidspunkt for lydhændelsen til operatøren, samt logger hændelsen. \newline
7. Use case afsluttes. 
\\ \hline
\textbf{Udvidelser/undtagelser:} & 

\textbf{2a. Lydhændelsen detekteres ikke af alle mikrofoner} \newline
1. Silent kasseres lydhændelsen \newline
2. Use case afsluttes. \newline
\newline
\textbf{5a. Systemet vurderer at estimatet ikke er troværdigt} \newline
1. Systemet viser den estimerede lokation til operatøren, men markerer den som utroværdig, med en fejlmargin vist. \newline
2. Use case fortsætter fra trin 6. \newline
\newline
\textbf{5b. Den estimerede lokation er udenfor måleområdet} \newline
1. Systemet markerer positionen som udenfor måleområdet. \newline
2. Use case fortsætter fra trin 6. \newline
\newline
\textbf{5b. En sensornode mister forbindelse, eller kommer} \newline \textbf{ud af sync på et vilkårligt tidspunkt} \newline
1. Systemet afbryder behandling af lydhændelsen. \newline
2. Systemet informerer operatøren om fejlen. \newline
3. Use casen afsluttes.

\\ \hline
\textbf{Datavariationsliste:} & x \\
\hline
\end{tabular}
\renewcommand{\arraystretch}{1}
\end{center}

\newpage
\subsection*{UC2 - Konfigurering af opstilling}

\begin{center}
\renewcommand{\arraystretch}{1.4}
\begin{tabular}{| p{150pt} | p{300pt} |}
\hline

\textbf{Navn:} & Konfigurér opstilling \\ \hline
\textbf{Mål:} & Operatøren kan indstille systemet med antal af mikrofoner samt deres lokation i måleområdet, under runtime af systemet.\\ \hline
\textbf{Aktører:} & Operatør \\ \hline
\textbf{Beskrivelse:} & Brugeren konfigurerer opstillingen af mikrofoner. Dette gøres via. terminalen, mens systemet kører. Operatøren indtaster antallet af mikrofoner, samt deres position (x,y).  \\ \hline




\hline
\end{tabular}
\renewcommand{\arraystretch}{1}
\end{center}

Navneord:
 Operatør
 Lydhændelse
 Mikrofon
 Lydhændelse
 Interface
 sensormoduler
 System-styring

 Verber:
 Interfacer
 Estimerer
 Detekterer
 Præsenteret (den estimerede lokation)
 Synkroniseret
 Vurderer
 
 
 
 


\end{document}